\uuid{eGnR}
\titre{Logique}
\auteur{bodin ; croizat ; sacre}
\organisation{exo7}
\licence{CC BY-NC-SA 4.0}
\datecreate{2022-03-08}
\niveau{L1}
\module{Algèbre}
\chapitre{Logique, ensemble, raisonnement}
\sousChapitre{Logique}
\difficulte{1}
\theme{Logique}
\exo7qcmid{hilisit.logique.001}

\contenu{
\begin{question}
\qclassification{100.01}
Quelles sont les assertions vraies ?
\begin{answers}
  \good{Il existe des triangles rectangles.}
  \bad{Tout triangle est un triangle rectangle.}
  \good{Tout triangle équilatéral est isocèle.}
  \bad{Il existe un triangle équilatéral qui est rectangle.}
\end{answers}
\begin{explanations}
Il existe des triangles rectangles, mais ce n'est pas le cas de tous les triangles.
Un triangle équilatéral (trois côtés égaux) est aussi isocèle (deux côtés égaux).
Un triangle équilatéral ne peut pas être un triangle rectangle (par le théorème de Pythagore ou parce que les angles d'un triangle équilatéral sont de $60^\circ$ tandis que l'angle droit mesure $90^\circ$).
\end{explanations}
\end{question}
}
